\documentclass[11pt,a4paper]{article}
\usepackage{templates/modeling-notes}
\hypersetup{pdftitle={Lecture 2 -- Linear programming: models and geometry}}
\title{\textbf{Mathematical Modeling}\\[2mm]
\large Lecture 2: Linear Programming---Models and Geometry}
\author{Study notes based on Ewa Pawelec's historical slides}
\date{}
\begin{document}
\maketitle
\noindent\textbf{Primary source:} \texttt{s4\_w2.pdf}, PDF pages 1--11,
printed slides 1--44.\\
\textbf{Source date:} Not established; stored in the folder labelled 2020.\\
\textbf{Current course alignment:} Unconfirmed.\\
\textbf{Editorial scope:} Source models and results are reorganized in English.
Hand derivations, proofs, cautions, and MATLAB checks are added study material.
\tableofcontents

\section{The linear programming problem}\label{sec:lp}
\textit{Source: slides 1--6, PDF pp.~1--2.}
A linear program (LP, Polish abbreviations PL or ZPL) optimizes a linear function
$c^Tx$ subject to linear equality and inequality constraints, with continuous
variables. The feasible set $\Omega$ includes every sign and bound restriction.
For maximization, a feasible $x^*$ is globally optimal when
$c^Tx^*\geq c^Tx$ for every $x\in\Omega$.

Minimization and maximization are related by
\[
 \min_{x\in\Omega}c^Tx=-\max_{x\in\Omega}(-c)^Tx.
\]
The optimizers agree after this sign change; the objective values have opposite
signs. An equivalent reformulation with additional variables requires a mapping
back to original decisions, not literal equality of vector sets in different dimensions.
Model construction follows Lecture 1: variables, constraints, objective,
algorithm, verification, and post-optimal analysis.

\section{Production planning: formulation and full hand solution}\label{sec:production}
\textit{Source: slides 7--9, PDF pp.~2--3; solution expanded.}
Two products consume three resources. The source data are:
\begin{center}
\begin{tabular}{lrrr}\toprule
 & Product 1 & Product 2 & Available\\\midrule
Resource 1 (kg/item; availability in kg) & 2 & 1 & 10\\
Resource 2 (kg/item; availability in kg) & 3 & 3 & 24\\
Resource 3 (kg/item; availability in kg) & 1.5 & 0 & 6\\
Selling price (PLN/item) & 3 & 2 & ---\\\bottomrule
\end{tabular}
\end{center}
Let $x_1,x_2$ be production quantities. Assume constant coefficients, additive
resource usage, and fractional production allowed. Since costs are absent,
$3x_1+2x_2$ is \emph{revenue}; the slides also call it profit, which would require
additional assumptions about costs. The model is
\begin{equation}\label{eq:production}
\max (3x_1+2x_2),\qquad
\begin{cases}
2x_1+x_2\leq10,\\
3x_1+3x_2\leq24,\\
1.5x_1\leq6,\\
x_1,x_2\geq0.
\end{cases}
\end{equation}
The second inequality is $x_1+x_2\leq8$, and the third is $x_1\leq4$.
Intersect boundary lines and retain only feasible intersections:
\begin{center}
\begin{tabular}{rrrr}\toprule
Vertex & $x_1$ & $x_2$ & $3x_1+2x_2$\\\midrule
$V_1$ & 0 & 0 & 0\\
$V_2$ & 4 & 0 & 12\\
$V_3$ & 4 & 2 & 16\\
$V_4$ & 2 & 6 & 18\\
$V_5$ & 0 & 8 & 16\\\bottomrule
\end{tabular}
\end{center}
For $V_4$, subtract $x_1+x_2=8$ from $2x_1+x_2=10$ to obtain $x_1=2$,
then $x_2=6$. These give resource slacks $(0,0,3)^T$ kg.

\begin{center}
\setlength{\unitlength}{5mm}
\begin{picture}(14,12)(-2,-2)
\put(0,0){\vector(1,0){11}}\put(0,0){\vector(0,1){10}}
\put(11,-0.7){$x_1$}\put(-1,9.5){$x_2$}
\put(0,0){\line(1,0){4}}\put(4,0){\line(0,1){2}}
\put(4,2){\line(-1,2){2}}\put(2,6){\line(-1,1){2}}
\put(0,8){\line(0,-1){8}}
\put(0,9){\line(2,-3){6}}
\put(2,6){\circle*{0.25}}\put(2.3,6.2){$(2,6)$}
\put(1,2){$\Omega$}\put(4.2,2){$(4,2)$}
\put(-1.5,8){$8$}\put(4,-0.8){$4$}
\put(6.2,0.7){$3x_1+2x_2=18$}
\end{picture}
\end{center}
The objective level lines have slope $-3/2$. Moving them toward increasing
revenue brings the last feasible contact to $(2,6)$.

\subsection{An optimality certificate without a graph}
For every feasible point,
\[
3x_1+2x_2=(2x_1+x_2)+(x_1+x_2)\leq10+8=18.
\]
The candidate $(2,6)$ attains the bound, proving global optimality. Equality
requires both resource inequalities to be tight, whose unique intersection is
$(2,6)$; thus this optimum is unique. This reasoning is more informative than
merely listing a solver output.

\section{Other modeling patterns}\label{sec:models}
\subsection{Diet problem}
\textit{Source: slides 10--11, PDF p.~3.}
Let $x_j$ be the daily quantity of food $j$, $c_j$ its unit price, and $a_{ij}$
the amount of nutrient $i$ per food unit. Nutrient limits are $b_i,d_i$, while
food-quantity bounds are $g_j,h_j$. The model is
\[
\min\sum_{j=1}^n c_jx_j,\qquad
b_i\leq\sum_{j=1}^n a_{ij}x_j\leq d_i\ (i=1,\ldots,m),\quad
g_j\leq x_j\leq h_j,\quad x_j\geq0.
\]
Nutrient and food bounds have different units and meanings. Positive food lower
bounds can enforce minimum variety, but generic bounds do not automatically
produce a realistic diet. No numerical instance is provided in these slides.

\subsection{Transportation problem}
\textit{Source: slides 12--13, PDF pp.~3--4.}
Let $x_{ij}$ be the quantity shipped from supplier $i$ to customer $j$;
$a_i$ is available supply, $b_j$ required demand, and $c_{ij}$ cost per shipped
unit. The source formulation is
\[
\min\sum_{i=1}^m\sum_{j=1}^n c_{ij}x_{ij},\qquad
\sum_{j=1}^n x_{ij}\leq a_i,\quad
\sum_{i=1}^m x_{ij}\geq b_j,\quad x_{ij}\geq0.
\]
\textbf{Added qualification:} With nonnegative supply/demand and every route
available without capacity limits, feasibility is equivalent to
$\sum_i a_i\geq\sum_j b_j$. Necessity follows by summing constraints;
sufficiency follows by allocating supplies successively to remaining demands.
With restricted routes, the total-supply test alone is insufficient.
If all supply must be shipped and every demand met exactly, use equalities and
balanced totals instead. These are different modeling requirements.

\subsection{Regression as a linear program}
\textit{Source: slide 14, PDF p.~4; formulas checked against the rendered slide.}
For observed pairs $(t_i,y_i)$, fit $y\approx at+b$, with residuals
$e_i=y_i-at_i-b$. The $t_i$ are given observations, not decision variables.
The slope $a$ and intercept $b$ are free real variables.
Minimizing the largest absolute residual introduces $r\geq0$:
\begin{align*}
\min_{a,b,r}\quad&r,\\
at_i+b-r&\leq y_i,\\
-at_i-b-r&\leq-y_i\qquad(i=1,\ldots,N).
\end{align*}
Indeed, the two inequalities are equivalent to $-r\leq e_i\leq r$.
For the alternative least-absolute-deviation criterion, introduce one $r_i\geq0$
per observation and minimize $\sum_i r_i$ under $-r_i\leq e_i\leq r_i$.
The slides mention least squares as a common regression method, but
$\min\sum_i e_i^2$ is a quadratic problem, not the LP displayed here.
The three error criteria need not produce the same fitted line.

\section{Cutting stock and integer decisions}\label{sec:cutting}
\textit{Source: slides 15--19, PDF pp.~4--5.}
Stock has length $L$. Required pieces have lengths $\ell_i$ and demand $b_i$.
A cutting pattern $j$ specifies nonnegative integer piece counts $a_{ij}$,
satisfying $\sum_i a_{ij}\ell_i\leq L$. Its offcut is
\[
r_j=L-\sum_i a_{ij}\ell_i.
\]
Let $x_j$ count how many stock pieces use pattern $j$.
For exact demand the source gives
$\min r^Tx$ subject to $Ax=b$, $x\geq0$; its following example uses $Ax\geq b$.
The latter permits overproduction. Physical counts require integer $x_j$;
real-valued $x_j$ define an LP relaxation.

The source condition $r_j<\min_i\ell_i$ describes \emph{maximal} patterns:
no additional piece fits. It is not a requirement for every feasible pattern.
Under exact demand, nonmaximal patterns may be needed to avoid excess pieces.

\subsection{Cabinet example with 5.2 m stock: full hand solution}
\textit{Source: slide 18; derivation added.}
At least 200 cabinet sets require 1400 short pieces (0.7 m) and 800 long pieces
(2.5 m). The three \emph{listed technological patterns} are:
\begin{center}
\begin{tabular}{lrrr}\toprule
 & Pattern 1 & Pattern 2 & Pattern 3\\\midrule
Short pieces & 7 & 3 & 0\\
Long pieces & 0 & 1 & 2\\
Offcut (m) & 0.3 & 0.6 & 0.2\\\bottomrule
\end{tabular}
\end{center}
We optimize over these supplied patterns; they are not an enumeration of all
length-feasible cuts. For instance, six short pieces and no long ones would be
a feasible but nonmaximal cut. Assume the listed patterns are the permitted ones,
ignore saw kerf, and allow surplus pieces as in the source model:
\[
\min 0.3x_1+0.6x_2+0.2x_3,\qquad
7x_1+3x_2\geq1400,\quad x_2+2x_3\geq800,\quad x\geq0.
\]
The candidate $(200,0,400)^T$ meets both demands exactly and has offcut
$60+80=140$ m. To prove optimality, take nonnegative multiples $3/70$ and $1/10$
of the two demand inequalities. Their left-hand side is
\[
\frac3{10}x_1+\frac8{35}x_2+\frac15x_3\geq140.
\]
Since $0.6-8/35=13/35>0$,
\[
0.3x_1+0.6x_2+0.2x_3\geq140+\frac{13}{35}x_2\geq140.
\]
Equality requires $x_2=0$ and both demands tight, yielding the stated candidate.
It is the unique optimum for the supplied continuous model and, being integer,
also solves its integer version. It uses 600 stock boards.

\begin{remark}[Offcut versus stock consumption]
With exact demand, $r^Tx=L\sum_jx_j-\sum_i\ell_i b_i$, so minimizing offcut
and minimizing stock count are equivalent. With minimum demand, produced length
can vary; the objectives are not generally equivalent. Account separately for
unwanted surplus, saw kerf, and stock limits when the real task requires them.
\end{remark}

\subsection{The 2 m cabinet prompt: an explicitly assumed interpretation}
\textit{Source: slide 19; solution added under the following assumptions.}
The original prompt is:
\begin{otherlanguage}{polish}
\begin{quote}
Ile desek o długości d=2 m potrzeba na wykonanie 7 szafek?
\end{quote}
\end{otherlanguage}
\textbf{English translation:} How many 2 m boards are needed to make seven cabinets?
The sketch shows three horizontal members and two vertical members. Interpret
the marked dimensions as cut lengths: three 0.8 m and two 0.3 m pieces per cabinet.
This assumes zero kerf and ignores board thickness and joint allowances; the
sketch does not specify those details. It gives demands 21 long and 14 short pieces.

Writing patterns as (long pieces, short pieces), the maximal patterns are
$(2,1)$, $(1,4)$, and $(0,6)$. Let their counts be $u,v,w\in\{0,1,2,\ldots\}$.
The model is
\[
\min u+v+w,\qquad 2u+v\geq21,\quad u+4v+6w\geq14.
\]
Required total length is $21(0.8)+14(0.3)=21$ m, so at least
$\lceil21/2\rceil=11$ boards are necessary. Ten $(2,1)$ boards and one $(1,4)$
board produce exactly $(21,14)$ pieces, using 11 boards and leaving 1 m of offcut.
This construction attains the length lower bound and proves optimality under the
stated interpretation. It is not a claim about unspecified manufacturing details.

\section{Standard form and equivalent transformations}\label{sec:standard}
\textit{Source: slides 20--24, PDF pp.~5--6.}
The lecture's convention is
\begin{equation}\label{eq:standard}
\max c^Tx,\qquad Ax=b,\quad x\geq0,
\quad A\in\R^{m\times n},\quad b\in\R^m,\quad c\in\R^n.
\end{equation}
Other texts use different ``standard form'' conventions. Here it means equality
constraints and nonnegative variables. Full row rank is a useful basis assumption,
not part of the definition of every LP. Remove redundant equalities after checking
consistency; $m<n$ is a common case, not a universal requirement.

\begin{center}
\begin{tabular}{ll}\toprule
Original condition & Replacement\\\midrule
$a^Tx\leq b$ & $a^Tx+s=b$, $s\geq0$ (slack)\\
$a^Tx\geq b$ & $a^Tx-s=b$, $s\geq0$ (surplus)\\
$\ell\leq a^Tx\leq u$ & Two inequalities, transformed separately\\
$x_j\in\R$ & $x_j=x_j^+-x_j^-$, $x_j^+,x_j^-\geq0$\\
$x_j\geq\ell_j$ & $x_j=\ell_j+y_j$, $y_j\geq0$\\\bottomrule
\end{tabular}
\end{center}
A finite upper bound becomes $y_j\leq u_j-\ell_j$ after the last substitution.
Substitution may add a constant to the objective; retain it when reporting the
original objective value. Free-variable splitting is not unique: adding the same
positive number to $x_j^+$ and $x_j^-$ leaves $x_j$ unchanged.

\subsection{The source conversion example}
\textit{Source: slide 24.}
The original problem is
\[
\min 2x_1-x_2+5x_3-2x_4,\qquad
\begin{cases}
x_1+x_2+2x_3+5x_4\leq20,\\
2x_1-x_2-x_4\geq15,\\
3x_2+x_3-2x_4=14,\\
x_1,x_3,x_4\geq0,\quad x_2\in\R.
\end{cases}
\]
Set $x_2=p-q$ with $p,q\geq0$, add slack $s$ to the first constraint and
subtract surplus $t$ from the second. For $y=(x_1,p,q,x_3,x_4,s,t)^T\geq0$,
\[
\max\begin{pmatrix}-2&1&-1&-5&2&0&0\end{pmatrix}y,
\quad
\begin{pmatrix}
1&1&-1&2&5&1&0\\
2&-1&1&0&-1&0&-1\\
0&3&-3&1&-2&0&0
\end{pmatrix}y=
\begin{pmatrix}20\\15\\14\end{pmatrix}.
\]
For a feasible original point choose $p=\max(x_2,0)$, $q=\max(-x_2,0)$,
$s=20-x_1-x_2-2x_3-5x_4$, and $t=2x_1-x_2-x_4-15$.
Conversely, any feasible $y$ recovers a feasible original point by $x_2=p-q$.
The transformed objective is the negative of the original one.

\subsection{Linear equations alone do not decide LP feasibility}
\textit{Source: slide 21; qualifications added.}
The system $Ax=b$ is consistent exactly when
$\operatorname{rank}(A)=\operatorname{rank}([A\mid b])$.
If consistent with rank $r$, its solution space has $n-r$ free dimensions:
uniqueness requires $r=n$. However, an LP also requires $x\geq0$.
For example, $x=-1$ is a consistent equation with no nonnegative solution.
A rank test is therefore necessary but insufficient for standard-form feasibility.
Equivalently, writing $a_j$ for the columns of $A$, feasibility means
$b=\sum_jx_ja_j$ with $x_j\geq0$: $b$ belongs to their conical hull
(source slide 34).

\section{Geometry of feasible sets}\label{sec:geometry}
\textit{Source: slides 30--35, PDF pp.~8--9. Definitions corrected and expanded.}
\begin{definition}[Convexity]
A set $P$ is convex if $\lambda u+(1-\lambda)v\in P$ for any $u,v\in P$ and
$0\leq\lambda\leq1$.
\end{definition}
For $a\neq0$, $a^Tx=b$ defines a hyperplane, with closed halfspaces
$a^Tx\leq b$ and $a^Tx\geq b$. A \emph{polyhedron} is the intersection of finitely
many closed halfspaces; equalities can be represented by two inequalities.
A bounded polyhedron is a \emph{polytope}. These English terms distinguish what
the source calls \emph{zbi\'or wielo\'scienny} and, in the bounded case,
\emph{wielo\'scian}.

\begin{proposition}
An LP feasible set is a closed convex polyhedron.
\end{proposition}
\begin{proof}
Each linear inequality defines a closed halfspace. If $a^Tu\leq b$ and
$a^Tv\leq b$, then $a^T(\lambda u+(1-\lambda)v)\leq b$.
Intersections preserve convexity and closedness, including paired inequalities
for equalities and halfspaces for sign restrictions.
\end{proof}

\begin{definition}[Extreme point]\label{def:vertex}
A point $v\in P$ is an extreme point (vertex) if
$v=\lambda u+(1-\lambda)w$, $u,w\in P$, $0<\lambda<1$, implies $u=w=v$.
\end{definition}
The strict restriction on $\lambda$ matters: taking $\lambda=0$ or $1$ cannot
force both points to equal $v$. A point in the middle of a feasible edge is not
a vertex even if it is optimal.

\subsection{Recession directions and extreme rays}
A nonzero vector $d$ is a recession direction if $x+\alpha d\in P$ for every
$x\in P$ and $\alpha\geq0$. For a nonempty polyhedron,
\[
P=\{x:Gx\leq h\}\ \Longrightarrow\ \operatorname{rec}(P)=\{d:Gd\leq0\},
\]
where the recession cone includes $d=0$. Substitution proves sufficiency;
if a row had $G_i d>0$, sufficiently large $\alpha$ would violate it, proving
necessity. Consequently,
\[
P=\{x:Ax=b,\ x\geq0\}\ \Longrightarrow\
\operatorname{rec}(P)=\{d:Ad=0,\ d\geq0\}.
\]
A nonempty polyhedron is unbounded exactly when its recession cone has a nonzero
vector. The forward implication follows from a finite polyhedral representation;
the reverse follows from the feasible ray. The source states this for arbitrary
convex sets, which needs qualification: the nonclosed convex set
$\{(0,0)\}\cup\{(u,v):u>0,\ 0<v<1\}$ is unbounded but has no nonzero recession
direction. Any ray from the origin either leaves $0<v<1$ eventually or stays
on an excluded part of its boundary. The equivalence holds for nonempty closed
convex sets, including polyhedra.

For a pointed cone, a nonzero $d$ generates an \emph{extreme ray} if
$d=u+v$ with $u,v$ in the cone forces both $u$ and $v$ to be nonnegative multiples
of $d$. A ray is a set of positive multiples, so its generating vector has no
unique scale. ``Pointed'' means the cone contains no line through zero.

\section{Existence, unboundedness, and the vertex principle}\label{sec:outcomes}
\textit{Source: slides 25--29 and 35--40, PDF pp.~7--10.}
For an LP there are three objective outcomes: infeasible; feasible with objective
unbounded in the improving direction; or feasible with a finite attained optimum.
The last case may have one or infinitely many optimal decisions.

\subsection{Unbounded feasible set versus unbounded objective}
If $x^0$ is feasible in standard form and $Ad=0$, $d\geq0$, $c^Td>0$, then
\[
x^0+\alpha d\in P,\qquad c^T(x^0+\alpha d)=c^Tx^0+\alpha c^Td\to+\infty.
\]
Feasibility is an essential assumption: a direction alone does not cure an empty
feasible set. For the source example $x_1-x_2=1$, $x\geq0$, maximizing $x_1$,
use $x^0=(1,0)^T$ and $d=(1,1)^T$. The objective is $1+\alpha$.
On the same feasible ray, maximizing $-x_1$ instead has the unique optimum
$(1,0)$, although the feasible set remains unbounded.

\subsection{Representation by vertices and rays}
\begin{theorem}[Polyhedral representation, stated without full proof]\label{thm:representation}
For a nonempty pointed polyhedron $P$, let $v^1,\ldots,v^k$ be its vertices and
$d^1,\ldots,d^\ell$ representatives of its extreme recession rays. Then
\[
x\in P\quad\Longleftrightarrow\quad
x=\sum_{i=1}^k\lambda_i v^i+\sum_{j=1}^{\ell}\mu_j d^j,
\qquad \lambda_i,\mu_j\geq0,\quad\sum_i\lambda_i=1.
\]
For bounded $P$, the ray sum is absent.
\end{theorem}
Standard-form feasible sets are pointed: a line contained in $x\geq0$ would
force its direction to be zero. Thus the theorem applies when they are nonempty.
The source labels its representation result as Carath\'eodory's theorem.
The vertex--ray representation is usually called a polyhedral representation
or resolution theorem; Carath\'eodory's convex-hull theorem concerns reducing a
convex combination to at most $n+1$ points. These are related but distinct results.
The pointedness hypothesis is checked against the supplementary MIT reference
in Section~\ref{sec:sources}; polyhedra with lines need a lineality term.

\begin{theorem}[Optimal vertex for the lecture's standard form]\label{thm:vertex}
Let $P=\{x:Ax=b,\ x\geq0\}\neq\varnothing$. If $c^Tx$ is bounded above on $P$,
its maximum is attained at a vertex of $P$.
\end{theorem}
\begin{proof}[Proof using Theorem~\ref{thm:representation}]
If any $c^Td^j>0$, its feasible ray makes the objective unbounded. Hence all ray
contributions are nonpositive, and every feasible $x$ satisfies
\[
c^Tx=\sum_i\lambda_i c^Tv^i+\sum_j\mu_j c^Td^j
\leq\sum_i\lambda_i c^Tv^i\leq\max_i c^Tv^i.
\]
A vertex attaining that finite maximum is feasible and optimal.
\end{proof}
This is an \emph{existence} statement: not every optimum must be a vertex.
It also cannot be applied without qualification to arbitrary feasible sets.
For example, maximize $-x_1$ on $\{x\in\R^2:x_1\geq0\}$.
The optimum is zero on an entire line, and that halfspace has no vertices.

\begin{proposition}[Multiple optima]
Every convex combination of optimal LP solutions is optimal.
\end{proposition}
\begin{proof}
For optimal $u,v$ with common value $z^*$, feasibility follows from convexity and
$c^T(\lambda u+(1-\lambda)v)=\lambda z^*+(1-\lambda)z^*=z^*$.
\end{proof}
In $x_1+x_2\leq1$, $x\geq0$, maximizing $x_1+2x_2$ gives $(0,1)$ uniquely;
maximizing $2x_1+x_2$ gives $(1,0)$ uniquely. Maximizing $2x_1+2x_2$ makes the
whole edge $x_1+x_2=1$ optimal, as in source slides 27--28.
The representation above also allows zero-objective recession directions:
optimal decisions can lie on unbounded optimal rays, not just between vertices.

\section{Bases and basic feasible solutions}\label{sec:bases}
\textit{Source: slides 41--43, PDF p.~11; characterization and proof expanded.}
Assume $A\in\R^{m\times n}$ has full row rank $m\leq n$.
Choose $m$ independent columns indexed by $B$; their square matrix $A_B$ is
nonsingular. With $N$ the remaining indices, define
\[
x_N=0,\qquad x_B=A_B^{-1}b.
\]
This is a \emph{basic solution}. It is \emph{basic feasible} (BFS) only if
$x_B\geq0$. Compute it by solving $A_Bx_B=b$ rather than forming an inverse.

\begin{theorem}\label{thm:bfs}
A feasible point of $\{x:Ax=b,x\geq0\}$ is a vertex exactly when it is a BFS.
\end{theorem}
\begin{proof}
For a BFS $x=\lambda u+(1-\lambda)v$, $0<\lambda<1$, nonnegativity and $x_N=0$
force $u_N=v_N=0$. Then $A_Bu_B=A_Bv_B=b$, so nonsingularity yields $u=v=x$.

Conversely, let $S=\{j:x_j>0\}$ for a vertex $x$. If the columns $A_S$ were
dependent, some nonzero $d$ supported on $S$ would satisfy $Ad=0$.
For small $\eps>0$, both $x+\eps d$ and $x-\eps d$ would remain nonnegative
and feasible, contradicting extremality. Thus $A_S$ is independent and can be
extended to a basis $B$ of $m$ columns. The other variables are zero, so $x$ is
its basic feasible solution.
\end{proof}
A BFS has at most $m$ positive coordinates. That count alone is insufficient:
the corresponding columns must be independent. For example,
\[
A=\begin{pmatrix}1&1&0\\0&0&1\end{pmatrix},\quad b=(2,0)^T,
\quad x=(1,1,0)^T
\]
has only $m=2$ positive coordinates but is the midpoint of two feasible points,
$(2,0,0)^T$ and $(0,2,0)^T$, not a vertex.
A BFS is \emph{degenerate} when at least one of its $m$ basic coordinates is zero;
different bases can then represent the same vertex.

\subsection{Production example in standard form}
Add three slacks to Equation~\ref{eq:production}. The augmented vector is
$y=(x_1,x_2,s_1,s_2,s_3)^T$, with
\[
\begin{pmatrix}2&1&1&0&0\\3&3&0&1&0\\1.5&0&0&0&1\end{pmatrix}y
=\begin{pmatrix}10\\24\\6\end{pmatrix},\qquad y\geq0.
\]
At the origin in production coordinates, $B=\{3,4,5\}$ is the slack basis.
At the optimum, $B=\{1,2,5\}$ yields $(x_1,x_2,s_3)=(2,6,3)$ and
$s_1=s_2=0$. The determinant of this basis matrix is $3\neq0$.
The graphical intuition of simplex is to move between adjacent basic feasible
solutions toward improving objective values. Pivot rules, degeneracy handling,
and termination arguments require later algorithmic material; Lecture 2 does
not yet supply a complete simplex implementation.

\section{MATLAB checks and execution status}\label{sec:matlab}
The shared script is \texttt{matlab/exercises/lecture-02-checks.m}.
It enumerates production vertices, checks resource slacks and the optimal basis,
verifies the 5.2 m cutting certificate, and enumerates the explicitly assumed
2 m cabinet problem. It uses base MATLAB operations without Optimization Toolbox.
The inputs match the hand derivations. Assertions use absolute tolerance
$10^{-10}$ for these small examples; this is not a universal numerical tolerance.

\textbf{Execution status:} Not run in MATLAB; MATLAB is unavailable in the current
environment. Expected results are derived above: production $(2,6)$ with value
18 PLN and slacks $(0,0,3)$ kg; 5.2 m cutting $(200,0,400)$ with offcut 140 m;
2 m cabinet count 11 with pattern counts $(10,1,0)$.
Do not treat these expected values as recorded MATLAB output.

\section{Common mistakes and revision questions}\label{sec:revision}
\begin{itemize}[leftmargin=1.6em]
\item Do not confuse unbounded feasible sets with objectives unbounded above.
\item Do not call every feasible intersection a vertex; discard intersections
violating other constraints and check independent active constraints.
\item Do not confuse slack and surplus signs or omit nonnegativity.
\item Do not replace integer board counts by arbitrary fractional values without
calling the resulting problem an LP relaxation.
\item Do not replace exact demand by minimum demand without considering surplus.
\item A feasible basic solution needs nonnegative basic coordinates; a nonsingular
basis alone does not guarantee feasibility.
\end{itemize}
For revision: derive the production certificate without the figure; transform a
free variable and a $\geq$ constraint; construct an infeasible LP whose equations
are consistent; and explain why an optimal edge does not contradict the vertex theorem.

\section{Source map and Polish terminology}\label{sec:sources}
\begin{center}
\begin{tabular}{lll}\toprule
PDF pages & Printed slides & Content\\\midrule
1--2 & 1--8 & General formulation, modeling, production\\
3--5 & 9--20 & Model examples and standard form\\
6 & 21--24 & Rank and equivalent transformations\\
7--8 & 25--32 & LP outcomes, halfspaces, vertices\\
9--10 & 33--40 & Recession directions and vertex principle\\
11 & 41--44 & Bases and simplex intuition\\\bottomrule
\end{tabular}
\end{center}
The main source is \texttt{s4\_w2.pdf}; added proofs and worked solutions are
editorial explanations based on its models. The regression formula and cabinet
sketch were inspected visually because PDF text extraction loses their notation.
Polish vocabulary is maintained in \texttt{docs/glossary.md}.

\textbf{Supplementary cross-check:}
\href{https://ocw.mit.edu/courses/6-251j-introduction-to-mathematical-programming-fall-2009/ab890d20b071105239fe0b9bcf599f57_MIT6_251JF09_lec10.pdf}{MIT OCW, 15.081J/6.251J, Lecture 10 (Fall 2009)},
PDF pp.~4 and 6, sections on recession cones, unbounded LPs, and the resolution
theorem. Used to check hypotheses of the vertex--ray representation;
this supplement is not current course coverage. Its intended external storage
location and download status are recorded in \texttt{docs/sources.md}.
\end{document}
